\section{Threat Model and Problem Formulation}
\label{sec:threat}

To establish the technical requirements for VulnRAG, we must first formalize the environment in which it operates and the specific retrieval challenges it addresses. This section defines our system model, the attacker’s objectives, and the research questions that guide our evaluation.

\subsection{System Model and Attacker Objectives}
We consider an autonomous security assessment framework operating within a heterogeneous enterprise network. The system is tasked with performing continuous vulnerability discovery and penetration testing.
\begin{itemize}
    \item \textbf{Environment}: A dynamic landscape of services, software versions, and network configurations, many of which contain undisclosed or recently disclosed vulnerabilities.
    \item \textbf{Knowledge Base}: The system has access to a continuously updated vulnerability knowledge base (e.g., the National Vulnerability Database), which publishes thousands of new Common Vulnerabilities and Exposures (CVEs) annually.
    \item \textbf{Attacker Objective}: The system (acting as a red-teaming agent) aims to identify the most operationally significant vulnerabilities to establish an initial foothold or achieve privilege escalation.
\end{itemize}

The primary constraint in this model is not "data accessibility," but \textbf{Intelligence Precision}. With over 30,000 new CVEs published each year, the agent faces significant information overload. The threat to autonomous reliability is the failure to distinguish between "theoretical" vulnerabilities and "operationally actionable" threats.

\subsection{Problem Formulation: The Retrieval Gap}
We formalize the core research challenge as the \textbf{Retrieval Gap} ($\mathcal{G}_R$). Let $Q$ be a security query and $K$ be the set of all known vulnerabilities. A standard RAG system returns a subset $K_Q \subset K$ based on semantic similarity $S(q, k)$. The Retrieval Gap exists when:
\begin{equation}
\mathcal{G}_R = K_{true} \setminus K_Q \neq \emptyset
\end{equation}
where $K_{true}$ is the set of vulnerabilities that are truly operationally relevant to the target environment (e.g., those part of an attack chain or having public exploits), but which lack sufficient semantic similarity to $Q$ to be retrieved by naive vector search.

VulnRAG aims to minimize $\mathcal{G}_R$ by incorporating structural relationship traversal and agentic self-correction into the retrieval loop.

\subsection{Research Questions}
Our study is guided by four primary research questions (RQs) that evaluate the effectiveness of VulnRAG in bridging the Retrieval Gap:
\begin{itemize}
    \item \textbf{RQ1 (Structural Discovery)}: To what extent does graph-based traversal of CVE relationship metadata improve the recall of operationally linked vulnerabilities compared to standard semantic vector search?
    \item \textbf{RQ2 (Agentic Resilience)}: Can a zero-shot agentic feedback loop (Evaluator-Reformulator) achieve search precision comparable to fine-tuned self-reflective RAG models?
    \item \textbf{RQ3 (Ranking Quality)}: How does the integration of security-specific domain metrics (CVSS, exploitability, recency) into the reranking function influence the "Intelligence Density" of results?
    \item \textbf{RQ4 (Operational Economics)}: What are the quantitative trade-offs between reasoning success rates and inference costs when using complexity-aware multi-tier routing?
\end{itemize}
